\section*{Goldsponsor}
\begin{center}
	\includegraphics[width=0.6\textwidth]{001_Wheregroup}
\end{center}
Die WhereGroup GmbH \& Co. KG ist ein mittelständischer Dienstleister und Lösungsanbieter im Bereich webbasierter
Geodateninfrastrukturen. Wir bieten Beratung, Konzeption, Entwicklung, Aufbau und Betrieb
dynamischer Kartenanwendungen im Intra- und Internet. Darüber hinaus gehört ein umfangreiches
Schulungs- und Workshop-Programm zu unserem Portfolio. Gegründet wurde das Unternehmen 2007 als
eine Fusion der Unternehmen CCGIS GbR (gegründet 1998), KARTA.GO GmbH (gegründet 2002) und
Geo-Consortium GbR (gegründet 2004) in Bonn. Bei der WhereGroup sind derzeit 30 Angestellte
unterschiedlicher Fachrichtungen am Hauptsitz in Bonn sowie in Freiburg und Berlin
beschäftigt.

Das Spektrum unserer Dienstleistungen reicht von Geoportalen und kartenbasierter
Datenverwaltung bis hin zu hochverfügbaren Anwendungen für die freie Wirtschaft und die öffentliche
Verwaltung.

In unseren Projekten setzen wir auf die Standards bzw. Empfehlungen des Open Geospatial
Consortiums (OGC), der INSPIRE"=Richtlinie sowie der GDI"=DE. Ihre Verwendung gewährleistet ein Maximum
an Interoperabilität und Flexibilität unserer Lösungen. Die Einhaltung hoher Sicherheitsstandards
ist für uns nicht zuletzt durch unsere Projekte mit Landes- und Bundesbehörden sowie Großkonzernen
eine Selbst\-ver\-ständ\-lichkeit.

Wir beraten absolut herstellerunabhängig und sind spezialisiert auf die Weiterentwicklung,
professionelle Anwendung und Integration offener Standards und bewährter
Open"=Source"=Technologien und freier Software. Dazu zählen neben unseren Projekten Mapbender\,3,
MetaDor\,2 und PostNAS u.\,a. GeoServer, MapServer, MapProxy, OpenLay\-ers, PostGIS, QGIS und
OpenStreetMap.

Über unser Schulungsinstitut, die FOSS Academy, bieten wir praxisorientierte
Schulungen zum Thema "`GIS mit Open-Source-Software"' an. Diese können sowohl von Einzelpersonen als
auch von Firmen, auf Wunsch auch als Inhouse-Schulungen, gebucht werden.
Zu unserer Überzeugung
gehört, dass wir uns aktiv in die Geoinformatik-Community einbringen. Es ist uns wichtig, an der
Diskussion und Weiterentwicklung von verschiedensten Open-Source-Lösungen mitzuwirken.

\enlargethispage{1.5\baselineskip}
Die WhereGroup ist bundesweit und international mit Hochschulen, Firmen und Verbänden vernetzt. Wir
verfügen über langjährige, persönliche Kontakte zu diversen Universitäten und Hochschulen im In- und
Ausland, zum FOSSGIS e.V., zur Open Source Geospatial Foundation (OSGeo), zum Open Geospatial
Consortium (OGC), sowie zu den Herstellern bzw. Maintainern der gängigsten Open-Source-Produkte im
Geo-Bereich.  Mehr zur WhereGroup unter www.wheregroup.com und www.foss-academy.com.
